\documentclass[11pt,a4paper]{article}
\usepackage[margin=24mm]{geometry}
\usepackage{amsmath,amssymb,booktabs,longtable,array,hyperref}
\hypersetup{hidelinks}
\newcommand{\Lc}{\textbf{(c)}\ }
\newcommand{\Ld}{\textbf{(d)}\ }
\newcommand{\scr}{\texttt{audit\_\allowbreak resonant\_\allowbreak channels.py}}
\newcolumntype{L}[1]{>{\raggedright\arraybackslash}p{#1}}
\emergencystretch=3em
\title{Independent audit of resonant Boussinesq channels:\\classical charge conservation, heavy instability,\\and conditional quantum conclusions}
\author{Quantum Boussinesq solitons project}
\date{4 November 2026}
\begin{document}
\maketitle
\begin{abstract}
\Lc Every heavy soliton has an exactly on-shell, opposite-moving light pair,
with matching $M,\Pi,P,E$ and the classical hierarchy. Independent Riccati
calculations through $H_8$ confirm the printed charge decompositions; a Darboux
generating function establishes the chord structure without the package's
invalid commuting-background argument. The exact Y family also gives a
localized growing mode of every interior heavy soliton: instability is already
classical. The real one-loop mismatch is reproduced conditionally, not promoted
to a physical decay width. Printed threshold-slope and singular-example numbers,
and the elementary factor in the Gamma fusion formula, fail independent checks.
\Ld An exact quantum S-matrix, quantum charge selection rules and decay widths
remain unestablished.
\end{abstract}

\section{Scope, conventions and reproducibility}
Labels have the repository meanings: (a) literature theorem, (b) literature
heuristic, (c) independent project derivation/computational evidence, (d)
unestablished physical interpretation. ``Confirmed'' certifies the stated
mathematical scope; ``qualified'' means conditional, restricted or incompletely
justified; ``refuted'' identifies a concrete discrepancy. Each new calculation
is in \texttt{scripts/}\scr. No package code is imported. No git operation,
existing-file edit, bulk phase-script run or new renormalisation choice was made.

\Lc Read: \texttt{AGENTS.md}, \texttt{README.md}, the complete target TeX and its
16-page PDF text, all three target scripts, the format of
\texttt{audit\_phase2i.tex}, the ranked next steps, and relevant controlling
passages of \texttt{errata.tex}, \texttt{audit\_smx.tex},
\texttt{audit\_defq.tex}, \texttt{audit\_det\_projected.tex} and
\texttt{audit\_good\_renorm.tex}. The last audit matters because even the real
one-loop dispersion depends on the finite $u^2$ subtraction. No external
reference, abstract or full text was newly read; no inherited literature
attribution is independently authenticated here. All mathematical findings
below are labelled (c), rather than claiming an unread literature proof.

\Lc We retain
\[
 u_{tt}=u_{xx}-u_{xxxx}-3(u^2)_{xx},\quad u=\varphi_x=2(\log\tau)_{xx},\quad
 \pi=\varphi_t,\qquad \Omega=\int d\pi\wedge d\varphi.
\]
The physical densities are $M:\ u$, $\Pi:\ \pi$, $P:\ -\pi u$ and
$E:\ (\pi^2+u^2+u_x^2)/2-u^3$. Direct integration agrees with
$M=2k$, $\Pi=-2kv$, $P=2k^3v/3$, $E=2k^3(5-4k^2)/15$, $k^2+v^2=1$.
For the bad equation the dispersion instead has $k=\sinh\beta$,
$v=\epsilon\cosh\beta$, $\omega^2=k^2+k^4$; no bad-equation quantum
prescription is transferred here. The inherited source dictionary is
$u_{\rm CZ}=-3\varphi$, $\mathcal L_{\rm CZ}=9\mathcal L$.
Scattering uses $s=\sigma_{\rm vac}$, $q_D=-s/(2\sqrt3)$; the spectral
parameter is $\eta\sigma_{\rm vac}$. These must not be interchanged.
The target uses only an abstract directed-edge holonomy, so it does not
reproduce the refuted signed scattering substitution in its computations.

\section{Verdict table: all labelled claim groups}
Repeated claims in the abstract and conclusion are covered by their numbered
source proposition and the rows below. The equation names are the source's
LaTeX labels. Every row has independent checks in \scr\ unless explicitly
identified as an unestablished physical claim (d).
\small
\begin{longtable}{@{}L{5.6cm}L{2.4cm}L{7.2cm}@{}}
\toprule Claim & Verdict & Evidence and precise scope\\\midrule\endhead
C01. One-body and KP dictionary; cross-ratio (Sec.~1, Prop.~2.1) & Confirmed (c) & Bilinear polynomial, all four directions, physical charge identities.\\
C02. Three resonance curves (Prop.~2.1.1) & Qualified & Curves confirmed (c); printed $\ell>0$ fails for left-moving overtaking. Equal-soliton diagonal excluded.\\
C03. KP triangle and projections (Prop.~2.1.2; abstract) & Qualified & Main chord formulas confirmed (c). First abstract formula for $\alpha_{13}$ is false; static $\ell=0$ is an endpoint exception.\\
C04. Positive Y solution, splitting/merging (Prop.~2.2.1) & Qualified & Y solution confirmed (c). Literal heavy-plus-light fusion into the other light is refuted by positive mass; shared-root projections survive.\\
C05. Areas and measure (Prop.~2.2.2) & Confirmed (c) & Independent slice integrals: singular fractions $2/9$, $2/3$; two-body measure zero does not restrict heavy one-body measure.\\
C06. $H_1,\ldots,H_8$ and potentials (Thm.~3.1.1--2) & Confirmed (c) & New spectral series plus direct Riccati/beta integrations; printed $H_7,H_8$ correct.\\
C07. All-order proof and spin gap (Thm.~3.1.3) & Qualified & Conclusion established by a different formal Darboux proof (c); printed commuting-background argument invalid. Spin-gap assertion confirmed.\\
C08. Classical telescoping, topology, $\Pi$ (Prop.~3.2.1) & Confirmed (c) & All-order generator and $M,\Pi,P,E$ integrated through Y collision. Original script integrates only $P,E$.\\
C09. Mode shifts and $\Delta F=0.38879244\ldots$ (Prop.~3.2.2) & Qualified & Number and mode telescoping confirmed (c) by regulated root-transmission traces; finite subtraction and heavy-saddle determinant conditional (d).\\
C10. Shifted quantum fusion curve (Remark~3.3) & Qualified & Local $E,P$ curve persists (c); simultaneous quantum charges and physical dispersion unresolved (d).\\
C11. $h\log h$, delays and Taylor logarithm (Prop.~4.1) & Qualified & Generic classical coefficients confirmed (c), away from cusps; physical resummation/frame interpretation conditional (d).\\
C12. Light concavity and soliton $1\to2$ protection (Thm.~4.3.1) & Confirmed (c) & Symbolic velocity differentiation; concavity plus positive mass. Subadditivity requires a defined total light momentum.\\
C13. Every-heavy embedding and Jacobian (Thm.~4.3.2) & Qualified & Exact all-heavy charge matching confirmed (c); $1/(\sqrt3 k)$ is one ordered head-on contribution, not the unrestricted displayed integral.\\
C14. On-shell slope value (\texttt{onshell\_slope}) & Refuted (c) & Symbolic formula correct; printed $+0.241078\ldots$ is actually $-0.2664334283\ldots$.\\
C15. WKB deficit (Prop.~4.4) & Qualified & Damped-sheet algebra confirmed (c); opposite sheet amplifies. Neither a physical-sheet selection nor inelastic probability derived (d).\\
C16. Singular numeric example (Prop.~4.4) & Refuted (c) & All three printed $A$, $|\Delta\ell|$, $|S|^2$ values disagree.\\
C17. Classical and phonon fusion (Prop.~4.5.1--2) & Confirmed (c) & Hirota factors/phases and independently reconstructed phonon transmission telescope; absolute phase constants still conditional (d).\\
C18. Gamma fusion obstruction (Prop.~4.5.3) & Qualified & Implemented candidate fails, with quoted phase (c); printed elementary-factor exponent is reversed. Physical bootstrap premise unproved (d).\\
C19. Target-momentum remainder (Prop.~4.5.3) & Qualified & $0.8755292547\ldots$ logarithmic part confirmed (c); not the full Darboux Taylor remainder or a physical no-go.\\
C20. Four regimes, width, $2\times2$ block and chronological failure (Sec.~5; remarks) & Qualified (d) & Classical domain classification survives (c). No quantum vertex, spectral normalization, complete channels or quantum charges computed.\\
C21. Singular two-phase breathers & Qualified & Elementary bound $A_{\rm br}<1/4$ confirms singularity within this ansatz (c); does not exclude quantum bound states (d).\\
C22. Script completeness and timings (Sec.~6) & Qualified & All three scripts pass in 21.87 s (c); several advertised checks are absent or only assumed identities.\\
\bottomrule
\end{longtable}
\normalsize

\section{Hierarchy charges: a separate derivation}
\paragraph{Riccati and scattering generating functions (C06--C08).}
\Lc Put $L_0=D^3-D/4$ and solve $\lambda^3-\lambda/4=z^3$ with
$\lambda=z+O(z^{-1})$. For the eigenfunction logarithmic derivative $w=\psi_x/\psi$,
\begin{equation}
 w_{xx}+3ww_x+w^3+qw+r=z^3,\qquad
 w=z+\sum_{n\ge1}w_n z^{-n}.
 \label{riccati}
\end{equation}
The integrated formal trace identity is
$H_n=-n\int(w_n-w_{n,\rm vac})\,dx$; the discarded local terms are total
$x$ derivatives. It can be obtained by differentiating the formal eigenvalue
problem with respect to $z$ and cyclically moving $D$ under the integrated
residue. In particular $w_1=-q/3$, $w_2=(q_x-r)/3$, giving $H_1=\int(q-q_0)/3$
and $H_2=2\int r/3$ directly. Our recursive scalar ODE computation is independent
of the package's fractional pseudodifferential multiplication.

\Lc For $g=e^{-p_a x}+e^{-p_b x}$, $f=g_x/g$ and $B=D-f$,
$BL_0=LB$ with
\begin{equation}
 q=-\tfrac14+3f_x,\qquad r=3ff_x+3f_{xx},\qquad
 \psi=(\lambda-f)e^{\lambda x}.
\end{equation}
The vanishing constant coefficient is the derivative of
$f_{xx}+3ff_x+f^3-f/4=\text{constant}$; all exponentials have the same
$L_0$ eigenvalue. The endpoint scattering identity gives
\begin{align}
 \int(w-\lambda)\,dx
 &=\log\frac{\lambda+p_b}{\lambda+p_a}
 =-\sum_{j\ge1}\frac{\Delta_j}{j\lambda^j},\qquad
 \Delta_j=(-p_b)^j-(-p_a)^j,\\
 H_n&=n[z^{-n}]\sum_{j\ge1}\frac{\Delta_j}{j\lambda(z)^j}.
 \label{generator}
\end{align}
Every coefficient is a polynomial endpoint difference. Thus the formal local
hierarchy has chord additivity to all orders, not just to order eight.
The common vacuum spectrum is essential; this is not a binomial expansion of
noncommuting differential operators.

\Lc Independent series extraction in \scr\ gives precisely
\begin{align*}
&H_1=\Delta_1,\quad H_2=\Delta_2,\quad H_3=H_6=0,\\
&H_4=\Delta_4-\tfrac13\Delta_2,\qquad H_5=\Delta_5-\tfrac5{72}\Delta_1,\\
&H_7=\Delta_7-\tfrac7{12}\Delta_5+\tfrac7{324}\Delta_1,\qquad
 H_8=\Delta_8-\tfrac7{36}\Delta_4+\tfrac{11}{324}\Delta_2.
\end{align*}
All printed potentials for $M,\Pi,P,E,H_7,H_8$ have zero exact remainder
modulo $p_a^2+p_ap_b+p_b^2-1/4$. The pure-$k^3$ ansatz in
$\operatorname{span}\{\Delta_1,\Delta_5\}$ gives an inconsistent exact linear
system. Independently, direct scalar recursion through $w_8$ on
$u=2k^2y/(1+y)^2$, $y=e^{kx}$, reduces the integrals to
\[
 \int_0^\infty\frac{y^j}{(1+y)^N}\,dy
 =\frac{j!(N-j-2)!}{(N-1)!}.
\]
All eight charges agree exactly at $(k,v)=(1/2,\sqrt3/2)$ and
$(97/103,-20\sqrt3/103)$, including a heavy left-mover. These are exact
rational/algebraic evaluations, not fitted coefficients.

\Lc The target's proposed proof by shifting $\widetilde L$ by $-D/4$ is
not valid: $[D,\widetilde L]=\widetilde q_xD+\widetilde r_x$ is nonzero on a
soliton. Its vague operator logarithm is not a derivation of the claimed
zero-background trace formula. Equation~\eqref{generator} supplies the missing
argument within the formal scalar hierarchy. No claim about arbitrary
quantum/nonlocal charges follows from this formal trace theorem.

\paragraph{Across the actual collision.}
\Lc Use the minus-sign Y solution, $g=\sum_i\exp(-q_ix+\sqrt3q_i^2t)$, so the
same Darboux transformation applies with three terms. At every finite $t$,
$f(-\infty)=-q_1$, $f(+\infty)=-q_3$. Therefore the entire integrated generating
function remains $\log[(\lambda+q_3)/(\lambda+q_1)]$ through the collision.
Our checker integrates its spatial density at $t=-5,1,7$ and two finite
spectral parameters; it does not merely subtract endpoint formulas from
themselves. Separate spatial quadratures of all four physical densities at
these times match the heavy values. This includes $M$ and $\Pi$, unlike the
original script, whose actual quadrature loop selects only density indices
2 and 3 ($P,E$). Hierarchy additivity at collision was not directly tested by
that original quadrature.

\section{Heavy-soliton kinematics and classical instability}
\paragraph{Exact on-shell map, not an isolated heavy parameter (C12--C14).}
\Lc For a right-moving heavy soliton $a\in(\pi/3,\pi/2)$, set
\[
 \alpha_1=2\pi/3-a,\quad\epsilon_1=+1,\qquad
 \alpha_2=a-\pi/3,\quad\epsilon_2=-1.
\]
Then $\alpha_1+\alpha_2=\pi/3$ and the two chords concatenate into the heavy
chord. Reversing all velocities covers the left-moving branch. All four
physical charges and the entire formal hierarchy agree exactly. The curve in
the two-particle parameter plane has measure zero, but its projection onto
the heavy parameter covers the whole open heavy branch. These are distinct
measure statements.

\Lc No listed topological charge forbids this classical process: the potential
endpoints are fixed along the exact solution, $M$ is continuous rather than an
integer kink number, and $\Pi$ is explicitly conserved. The integrated spectral
generator also telescopes, so the usual local spectral hierarchy supplies no
additional obstruction. A smooth exact trajectory is stronger evidence than
matching four asymptotic numbers. \Ld This does not enumerate all possible
quantum superselection rules or compute an anomalous/nonlocal quantum charge.
The defensible quantum conclusion is compatibility with the checked classical
charges, not a theorem that no hidden quantum selection rule exists.

\Lc Using velocity as independent coordinate gives
\[
 \frac{dE}{dP}=v,\quad
 \frac{dP}{dv}=\frac23\sqrt{1-v^2}(1-4v^2),\quad
 \frac{d^2E_{\rm light}}{dP^2}=-\frac3{2\sin3\alpha}<0.
\]
Concavity and monotonicity imply the stated strict inequality when
$|P_1+P_2|\le P_{\max}$. If the sum lies outside this interval the right-hand
light energy is undefined, so the abstract's unrestricted wording needs this
restriction. Mass positivity excludes a heavy daughter in a light-soliton
$1\to2$ decay. This is soliton-sector kinematic protection, not a general
proof of quantum elastic factorisation with arbitrary radiation channels.

\Lc At fixed heavy momentum, the derivative of the head-on pair energy with
respect to its right-moving momentum is $v_1-v_2=\sqrt3 k_3>0$. Hence the
ordered head-on delta-function contribution is exactly
\begin{equation}
 J_{+-}=\frac1{\sqrt3 k_3}.
\end{equation}
It is finite and nonzero in the interior, with the strict range
$(1/\sqrt3,2/3)$; the printed closed lower endpoint requires the excluded
static limit. The displayed unrestricted integral
in the package needs a channel restriction (and a convention for identical
particle counting). If both orderings are integrated, both roots contribute.
There can also be co-moving $E,P$ roots near the cusp; our independent solver
finds one at heavy $a=1.1$. Its mass fails to match, so it does not undermine
the charge-conserving head-on result, but it refutes interpreting the displayed
integral as an unrestricted total phase space. Its widths are
$k_1=0.07962397547781387719\ldots$,
$k_2=0.83585099579192889532\ldots$, with mass mismatch
$M_1+M_2-M_3=0.04853522241661486512\ldots$.

\Lc The printed formula
$\Delta E'(0)=-(v_2-v_3)P_2'(\alpha_2)$ is correct. At $\alpha_1=0.62$ the
signed value is
\begin{equation}
 \Delta E'(0)=-0.266433428334333243278532392316\ldots,
\end{equation}
not $+0.24107867527655644$. We differentiate the energy mismatch using an
independent velocity inversion of the heavy momentum. The sub-threshold
statement $\Delta E(h)<0$ for $h>0$ survives; the printed positive
slope would have contradicted it. Globally in the light-light square,
$\partial_{\alpha_2}\Delta E=(v_2-v_3)P_2'(\alpha_2)>0$, since
$v_2<-1/2$ and $|v_3|<1/2$. The unique zero is the displayed head-on
resonance, so the lower triangle is below it. Independent sample checks
cover both sides.

\paragraph{An explicit classical unstable mode, absent from the package.}
\Lc In the plus-sign Y family insert a small coefficient $c_2=\varepsilon$.
In the frame $\xi=x-v_3t$, after removing a harmless exponential,
\begin{equation}
 \tau\ \sim\ 2\cosh(k_3\xi/2)
       +\varepsilon e^{-\sqrt3v_3\xi/2}e^{\gamma t},\qquad
 \gamma=\sqrt3 k_1k_2=\sqrt3(1/4-v_3^2)>0.
\end{equation}
Consequently $u=u_0+\varepsilon e^{\gamma t}\chi+O(\varepsilon^2)$ with
\begin{equation}
 u_0=\frac{k_3^2}{2}\operatorname{sech}^2(k_3\xi/2),\qquad
 \chi=2\partial_\xi^2
 \left[\frac{e^{-\sqrt3v_3\xi/2}}{2\cosh(k_3\xi/2)}\right].
 \label{unstable}
\end{equation}
Both tails decay precisely because $k_3>\sqrt3|v_3|$, equivalent to
$|v_3|<1/2$. Differentiating the exact family proves, and 80-digit numerical
spatial differentiation independently verifies,
\[
 (\gamma-v_3\partial_\xi)^2\chi
 =\chi''-\chi''''-6(u_0\chi)''.
\]
The associated potential perturbation also decays; its $M$ and $\Pi$
variations are zero boundary terms. Thus the heavy soliton already has a
localized exponentially growing classical mode. \Ld A real continuum mode
trace on this saddle is not by itself its physical quantum energy; discrete
unstable modes and the determinant contour/domain must be addressed. The
existence of this classical mode does not fix a quantum width or its
normalization.

\section{One-loop mismatch and direction of the shifted resonance}
\paragraph{Independent regulated computation (C09--C10).}
\Lc We reconstruct the phonon transmission from the soliton chord and complex
phonon roots
\[
 z_\pm(q)=\frac{\operatorname{sgn}(q)\sqrt{1+q^2}\pm i\sqrt3q}{2\sqrt3},\qquad
 A_{\rm ph}(q)=\frac{(p_a-z_+)(p_b-z_-)}{(p_a-z_-)(p_b-z_+)}.
\]
We differentiate its logarithm analytically and integrate the regulated
energy and momentum traces, including the inherited equal-mode endpoint and
tadpole:
\begin{align}
 e_\Lambda&=\frac1{4\pi}\int_{-\Lambda}^{\Lambda}
 |q|\sqrt{1+q^2}\,\theta'(q)\,dq-B_\Lambda
 +\frac{3k}{\pi}(\sqrt{1+\Lambda^2}-1),\\
 p_\Lambda&=\frac1{4\pi}\int_{-\Lambda}^{\Lambda}q\theta'(q)\,dq,\quad
 B_\Lambda=\frac{\Lambda\sqrt{1+\Lambda^2}}{4\pi}
 [\theta(\Lambda)-\theta(-\Lambda)].
\end{align}
Endpoint arguments are chosen to tend to zero at large $|q|$. No package
angular-log charge routine or saved result is used. With $q=\sinh t$,
independent cutoffs and a free intercept in $1/\Lambda,1/\Lambda^3$, the real
trace reproduces
\begin{equation}
 \Delta F=F_{12}+F_{23}-F_{13}
 =0.38879244070035460274\ldots,\qquad F=e-vp,
 \label{df}
\end{equation}
while $\Delta e=\Delta p=0$. This is conditional on that prescribed trace and
finite subtraction. The defect holonomy identity
$\sum_i\rho_i\mathcal O_i=\eta(\mathcal O_{12}+\mathcal O_{23}-\mathcal O_{13})$
is valid orientation bookkeeping, not a new physical absorption test. It is not a newly selected physical renormalisation.

\paragraph{What the positive sign means.}
\Lc At the old classical daughter momenta, keeping all three physical momenta
fixed, the corrected daughter energy exceeds the heavy energy by
$\hbar\Delta F>0$. Thus that particular old pair is off shell. It does not
follow that the entire decay closes: the classical root is inside the
head-on continuum and has a nonzero derivative. At fixed physical heavy $P_3$,
write $P_1\mapsto P_1+\hbar\delta p$, $P_2\mapsto P_2-\hbar\delta p$. Then
\begin{equation}
 \delta p=-\frac{\Delta F}{v_1-v_2}<0.
\end{equation}
For sufficiently small $\hbar$ at any fixed interior heavy parameter,
the right-moving daughter momentum decreases and the left-moving daughter
momentum becomes less negative. At $\phi=0.23$,
$\delta p=-0.23450418388533775558\ldots$. Both daughter widths decrease locally.
The equality moves toward $\alpha_1+\alpha_2<\pi/3$, the regular side, in
these fixed-physical-momentum coordinates. This is a shifted on-shell locus,
not a continuum-edge threshold opening/closing. The statement is nonuniform
near the heavy cusp, where one daughter becomes soft.

\Lc A finite term $\hbar\lambda u^2$, explicitly left unselected by the
controlling good-renormalisation audit, changes the mismatch by
\begin{equation}
 \Delta F_\lambda=\Delta F_0-2\lambda k_1k_2k_3.
\end{equation}
This follows from $\int u^2dx=2k^3/3$ and $k_3=k_1+k_2$, and is independently
checked. Therefore even the quoted numerical sign is a prescription-dependent
real-trace statement. \Ld Simultaneous conservation of the physical quantum
$M,\Pi,H_n$, and a physical complex energy of the unstable saddle, cannot be
inferred by shifting two momenta in an $E,P$ equation alone. The package
correctly acknowledges this higher-charge gap in its remark but overstates
instability/elasticity as settled in parts of the conclusion.

\section{Resonance phases, formal deficits and fusion}
\paragraph{Classical geometry and transverse limits (C01--C05, C11).}
\Lc Factoring the pairwise Hirota dispersion gives
\[
 \mathcal P(q_i-q_j,\sqrt3(q_i^2-q_j^2))
 =4(q_i-q_j)^2(q_i^2+q_iq_j+q_j^2-1/4).
\]
It proves the positive Y solution and the three shared-root loci. The correct
lines are $R_{\rm ho}^+:\alpha_1+\alpha_2=\pi/3$,
$R_{\rm ho}^-:|\alpha_1-\alpha_2|=\pi/3$, and
$R_{\rm ov}:\alpha_1+\alpha_2=2\pi/3$. No overtaking pole occurs in the open
square. The printed $\ell>0$ on overtaking is true only for right-movers;
left-movers have $\ell<0$. The first abstract expression
$\alpha_{13}=\alpha_{12}+\alpha_{23}+\pi/3$ would give $2\pi/3$,
outside the declared range, since the two light angles sum to $\pi/3$.
The main-text piecewise heavy angle is correct. At $\phi=\pi/6$ the cubic
is nondegenerate but the heavy chord is static, $k=1$, an excluded endpoint
of the stated open one-body parametrisation.

\Lc For generic noncusp points, differentiating the independently computed
Hirota phase twice gives
$\lim_{h\to0^+}h\delta''(h)=\frac29\sigma c_\ell$ with the three tabulated
coefficients. We use $h=10^{-25}$ and an independent derivative extraction,
not the package's two-point subtraction. The fixed-$P$ logarithmic coefficient
is nonzero at the three sampled points. The singular expansion is not uniform
when $h$ is comparable to $\hbar$; interpreting it as a resummed physical
branch displacement requires the missing physical frame/determinant.

\paragraph{Damping is a sheet-dependent candidate calculation (C15--C16).}
\Lc The two continuations $\log A=\log|A|\pm i\pi$ give reciprocal squared
moduli. On the damped branch only,
\[
 |e^{i\delta/\hbar}|^2=e^{-4\pi|\ell_1-\ell_2|/(9\hbar)}.
\]
Calling that branch ``physical'' because it is bounded is an assumption,
not a derivation of scattering boundary conditions. A deficient candidate
alone does not calculate the channels into which probability flows.
Away from the exact on-shell curve, $2\to1$ is not energy conserving, so it
cannot simply be assigned the entire open-domain deficit. The correct
numbers for the printed example $(0.50,0.65)$, directions $(+,-)$,
$\hbar=0.1$, are
\begin{align*}
 A&=-3.7826223332392473965427362479\ldots,\\
 |\ell_1-\ell_2|&=0.299443629683584018367638607696\ldots,\\
 |S|^2&=0.0152828833874182593860493271841\ldots.
\end{align*}
All three quoted package numbers fail; its script asserts only $A<0$ and
$0<|S|^2<1$, never those numbers. \Ld Identifying the complement with a
physical inelastic probability or proving failure of elastic quantum
scattering would require a controlled amplitude.

\paragraph{Fusion survives classically, but the Gamma candidate does not (C17--C19).}
\Lc Independent Hirota factors satisfy
$A(3,12)A(3,23)=A(3,13)$; the topological term telescopes using both $M$ and
$\Pi$. The full real classical phases therefore add. The independently
constructed complex phonon transmission itself telescopes pointwise, a
stronger test than subtracting angular phase derivatives. With the same
integration and reference prescription this proves additivity of the bare
integral linear in the target phase derivative. It does not set absolute
one-loop phase constants, precisely as the earlier phase audits warn.

\Lc Write $R(\eta)=\Gamma(1+i\eta)/\Gamma(1-i\eta)$. The implemented
Gamma candidate gives the fusion ratio
\begin{equation}
 \left[R(\eta)\left(\frac{e^2}{\eta^2}\right)^{+i\eta}
 \right]^{-\tau(3,13)},
 \label{gammafix}
\end{equation}
with a \emph{plus} exponent inside the brackets. The package prints a minus
exponent. The plus sign cancels the leading Stirling phase and tends to
$\exp[-i\pi\tau\operatorname{sgn}(\eta)/2]$. Fresh log-Gamma evaluations give
the quoted $1.5158110455271881045\ldots$ phase at $\hbar=0.1$ and approach
$+\pi/2$ at $10^{-4}$ and $10^{-8}$; the literal printed formula instead has phase
$2.88667602677875534274\ldots$ at $\hbar=0.1$.
Thus the obstruction for this candidate survives, while its displayed
identity is refuted. The earlier candidate audit does not identify this
candidate as the physical S-matrix.

\Lc Using the independently integrated mode momenta gives the target
logarithmic remainder $0.8755292547543045976\ldots$. Its nonvanishing is a
valid statement about the displayed expression. However the full Darboux
shift also contains $\gamma_jM_i$, so calling this logarithmic piece the
entire fixed-physical-momentum fusion mismatch is incomplete. In the same
mode-charge prescription, retaining that term gives
$3.58464625500292908012\ldots$ at the quoted point, rather than
$0.8755292547\ldots$. Moreover the
single-soliton momentum choice, two-body frame and finite subtraction remain
conditional. \Ld The proposed chronological failure and a universal physical
fusion no-go therefore remain conjectural.

\Lc The good two-phase breather bound follows directly from
$A_{\rm br}=(1-3\cot^2a)/(1+3\coth^2b)<1/4$. The tau polynomial cannot
remain positive when its interference minimum is reached with $A<1$;
this is an ansatz-specific classical singularity. \Ld It supplies neither a
quantum bound-state exclusion nor an exact S-matrix pole theorem.

\paragraph{Quantum claims that remain guesses (C20).}
\Ld The displayed decay-width normalization needs normalized states and a
defined vertex. The proposed $2\times2$ block assumes a complete two-channel
space, while an unstable heavy state is not automatically an asymptotic
one-particle channel. One cannot simultaneously infer its instability and
use it as an exact stable external state without specifying a resonance
formalism. Neither kinematic protection nor a scalar Yang--Baxter identity
proves factorisation. The four regimes are useful classical geometry, but
``must augment by $2\leftrightarrow1$'' throughout $A<0$ and ``only these
sectors can be elastic'' go beyond the checked results. The unresolved
quantum vertex and higher-charge corrections are substantive missing work.

\section{Verification coverage, numerical budgets and failures recorded}
\paragraph{Original scripts.}
\Lc Byte-identical copies were executed under new
\texttt{/tmp/audit\_\allowbreak resonant\_\allowbreak channels\_\allowbreak pkg\_*} names. Their only output
writes follow \texttt{\_\_file\_\_}, so their JSON files also stayed in
\texttt{/tmp}; tracked caches were untouched. Measured wall times were:
\begin{center}
\begin{tabular}{lrr}\toprule
Target script suffix & Seconds & Exit / final marker\\\midrule
\texttt{conservation.py} & 16.388 & 0 / present\\
\texttt{kinematics.py} & 3.709 & 0 / present\\
\texttt{quantum.py} & 1.773 & 0 / present\\\bottomrule
\end{tabular}
\end{center}
All print \texttt{ALL CHECKS PASSED}; the total is 21.870 s on this run,
consistent with the claimed under-25-second run, not a portable speed bound.

\Lc The hierarchy integrator is substantive independent-of-target code within
the package; it is not a fitted formula compared with itself. Conversely:
\begin{itemize}
\item The symbolic telescoping loop converts stringified expressions with
\texttt{sympify} without preserving the original real-symbol assumptions.
Its $p$ can differ from the substitution symbol, so these particular
substitutions can yield identically zero charges before the check.
Even without that bug, cancellation of a supplied $F$ with itself would
check bookkeeping rather than derivation. The earlier exact residue
integration does supply substantive independent evidence.
\item The claimed direct collision tests of $M,\Pi$ and all $H_n$ are absent;
only $P,E$ are spatially integrated. The all-order proof is asserted in prose.
\item The phase-space Jacobian is not numerically integrated or explicitly
asserted in the quantum checker. Area fractions are saved as strings rather
than verified by area integration. Bare phonon derivatives are sampled at
four points; the advertised continuum fusion integral is not integrated.
\item The unitarity test inserts the claimed positive imaginary part and
checks an exponential is between zero and one; it selects no physical sheet.
The printed slope and singular example are not compared with computed values.
The Gamma comparison uses the corrected plus sign, not the printed minus.
\item The one-loop mismatch test asserts only a nonzero discrepancy, not the
20-digit printed number; the independently reconstructed trace here supplies
that missing numerical cross-check. No quantum matrix element is tested.
\end{itemize}
\Lc No evidence of tolerance weakening in the available final scripts was
found. Their own failures list explicitly admits an initially fabricated
trailing-digit target, an incorrect early coefficient guess, insufficient
asymptotic time, and root-inversion/sign mistakes. Those admissions are not
erased by current passing tests. Editing history cannot be reconstructed
from these snapshots.

\paragraph{Independent error accounting.}
\Lc Exact symbolic identities have no mesh/fit error. Riccati beta integrals
are exact at the two stated profiles; the generic generating-function
identities are polynomial identities. Y quadratures use central moments,
different panels and times from the package, and 45/65-digit arithmetic;
finite-spectral-parameter generator integrals use 80 digits. Their differences
are numerical diagnostics, not rigorous interval error bounds. The root-trace
calculation separates (i) quadrature precision/panel refinement, (ii) finite
cutoff variation, and (iii) two- versus three-cutoff extrapolation spread;
no leading target coefficient or mismatch is fed into the fit. These budgets
certify conditional integrals, not the uncomputed determinant or a physical
renormalisation uncertainty. Asymptotic $h\log h$ extraction uses a declared
$10^{-21}$ tolerance at $h=10^{-25}$ rather than claiming all 80 digits.

\Lc The measured diagnostics are:
\begin{center}\small
\begin{tabular}{lr}\toprule
Quantity / independent variation & Value\\\midrule
$\Delta F$, $\Lambda=\sinh18$ & $0.38879243375721738313$\\
$\Delta F$, $\Lambda=\sinh24$ & $0.38879244068314428624$\\
$\Delta F$, $\Lambda=\sinh30$ & $0.38879244070031194263$\\
Free three-cutoff intercept & $0.38879244070035460274$\\
Two-/three-cutoff intercept spread & $6.10\times10^{-35}$\\
Trace quadrature method/precision/panel spread & $1.66\times10^{-47}$\\
Y physical-charge 45-/65-digit spread & $8.55\times10^{-47}$\\\bottomrule
\end{tabular}
\end{center}
The trace mesh check compares 80-digit tanh--sinh on $t$ panels
$0,1,3,6,18$ with 60-digit Gauss--Legendre on
$0,1/2,2,5,10,18$. The finite-cutoff discrepancy at the finest cutoff is
about $4.27\times10^{-14}$, much larger than quadrature error; it is removed
by extrapolation, not silently included in a high-precision equality.
The two-/three-cutoff spread is model sensitivity, not a statistical confidence
interval or an all-orders tail bound. No empirical fit is used for $H_n$.

\paragraph{Failures recorded.}
\begin{enumerate}
\item Package-claim discrepancies: left-moving overtaking $\ell$ sign;
first abstract heavy-angle formula; invalid commuting-background proof;
unrestricted phase-space formula; slope sign and magnitude; three singular
example numbers; Gamma elementary-factor sign; literal heavy-plus-light fusion labels. Refuted claims are encoded as
assertions that the discrepancies exist.
\item The first independent symbolic integration attempt grew too slowly
with rational $(k,v)$ carrying irrational coefficient denominators and was
interrupted. We changed exact test profiles so $k$ and $v/\sqrt3$ are rational,
retaining both a light and a heavy case and all eight charges. The generic
symbolic generating-function identity was retained, not weakened.
\item An initial audit check assumed that merely reversing the printed
slope sign would repair it. It failed with a $0.025354753\ldots$ magnitude
discrepancy. We replaced that false target with an independently differentiated
analytic expression and explicitly assert the printed magnitude discrepancy.
\item The first PDF build reported an overfull temporary-path line.
Explicit filename break points repaired it; all pages were rendered again.
\end{enumerate}

\paragraph{Measured independent execution.}
\Lc The final complete checker ran in 171.773 seconds and printed
\texttt{ALL CHECKS PASSED}. Earlier complete successful versions took
141.592 and 180.752 seconds; the final run includes the additional printed-table,
phase-space, breather and literal-fusion discrepancy assertions. All three
original script copies also printed the marker. Assertions encode refutations
as detected discrepancies, not false equalities; no failed tolerance was relaxed.
The final PDF was built twice, without overfull-box warnings, and all ten pages
were rendered and visually inspected.

\paragraph{Deliverables and references.}
The independent checker, this TeX and its twice-built PDF are the only new
repository files. No existing status or inventory file was edited, as required
by the task. Temporary build and inspection files have the same audit prefix.
No bibliography is invented: cited local audit filenames identify the exact
scope of inherited qualifications, and no external full-text reading is
claimed. \Ld The remaining scientific task is a quantum treatment of the
unstable heavy saddle and its coupled channels with physical higher charges,
not further refinement of already exact classical charge matching.
\end{document}
